% Source: https://artofproblemsolving.com/community/c3h2742823_gcd_property
% Stable identity: aops-topic-2742823-q01
Let $S=\{ax+by: x,y\in\mathbb Z\}$, and let $s$ be the smallest positive element of $S$ (it exists, since $a\in S$).

\emph{$s$ divides $a$ and $b$.} Write $a=qs+r$ with $0\le r<s$. Then $r=a-qs$ is in $S$, and since $r<s$ it cannot be positive, so $r=0$ and $s\mid a$. Similarly $s\mid b$.

\emph{$s=\gcd(a,b)$.} As a common divisor, $s\le d=\gcd(a,b)$. Conversely $d$ divides $a$ and $b$, hence divides $s=ax+by$, so $d\le s$. Thus $s=d$.

\emph{Conclusion.} Writing $d=ax+by$, any $c$ dividing $a$ and $b$ divides $ax+by=d$. $\blacksquare$

(This is why many books define the gcd by the divisibility property: in other number systems, such as polynomial rings, ``largest'' may not make sense, while ``divisible by every common divisor'' does.)
